%==============================================================================
% Sjabloon projectvoorstel graduaatsproject
%==============================================================================
% Je gegevens staan in ../projectgegevens.tex. De inhoud van je voorstel
% schrijf je in voorstel-inhoud.tex - dat bestand wordt later ook automatisch
% als bijlage in je rapport opgenomen.

\documentclass{hogent-article}

% Voor een voorstel in het Engels: \documentclass[english]{hogent-article}
% Dat kan enkel na toestemming van je begeleider.

\addbibresource{voorstel.bib}

%%=============================================================================
%% Projectgegevens - VUL DIT ALS EERSTE IN
%%=============================================================================
%%
%% Dit is de enige plek waar je je naam, je titel, je bedrijf en je begeleiders
%% invult. Het voorstel, het rapport en de poster halen hun gegevens hier op,
%% zodat ze nooit uit elkaar kunnen lopen.
%%
%% Enkel de tekst tussen de accolades aanpassen. Laat de commandonamen staan.

%% ---------- Jij ------------------------------------------------------------
\newcommand{\studentnaam}{Kelly De Clercq}
\newcommand{\studentemail}{kelly.declercq@student.hogent.be}

%% Wordt gebruikt als bestandsnaam van je PDF's, bijvoorbeeld
%% Familienaam-Voornaam-rapport.pdf. Geen spaties of leestekens gebruiken.
\newcommand{\bestandsnaam}{kelly-declercq}

%% ---------- Je project -----------------------------------------------------
\newcommand{\projecttitel}{perimeter-punches}

%% Een ondertitel is optioneel. Laat leeg als je er geen hebt: {}
\newcommand{\projectondertitel}{Optionele ondertitel}

%% Een of twee zinnen. Wordt gebruikt op de poster en in het voorstel.
\newcommand{\projectpitch}{%
  Een beveiligde applicatie die het wegvallen van terreininzicht na de migratie naar UKG opvangt, 
  door leidinggevenden real-time aanwezigheden aan de perimeter en historische tiktijden tot zes maanden te laten raadplegen. %
}

%% ---------- Waar je het doet ------------------------------------------------
\newcommand{\bedrijfsnaam}{Volvo Cars Gent}

%% Je mentor op de werkvloer (het bedrijf)
\newcommand{\mentornaam}{Tomas Bolckmans}
\newcommand{\mentoremail}{tomas.bolckmans@volvocars.com}

%% Je begeleider van de opleiding (HOGENT)
\newcommand{\begeleidernaam}{Luc Vervoort}
\newcommand{\begeleideremail}{luc.vervoort@hogent.be}

%% ---------- Je repository ---------------------------------------------------
%% De link naar je eigen git-repository. Die hoort zichtbaar in je rapport en
%% op je poster: hij is een beoordeeld onderdeel van je oplevering.
\newcommand{\repolink}{https://github.com/kellydeclercq/mijn-graduaatsproject.git} 

%% ---------- Vast voor dit opleidingsonderdeel -------------------------------
%% Deze twee hoef je normaal niet aan te passen.
\newcommand{\opleidingnaam}{Graduaat in het Programmeren}
\newcommand{\academiejaarwaarde}{2026-2027}

%% Zoekwoorden voor het voorstel en de poster (3 tot 5)
\newcommand{\projectkeywords}{tijdsregistratie, aanwezigheidsopvolging, toegangsbeveiliging}


%% Opleiding, vak en soort opdracht
\studyprogramme{\opleidingnaam}
\course{Graduaatsproject}
\assignmenttype{Projectvoorstel}

\academicyear{\academiejaarwaarde}

\title{\projecttitel}

\author{\studentnaam}
\email{\studentemail}

%% Je mentor op de werkvloer
\supervisor[Mentor]{\mentornaam{} (\bedrijfsnaam, \href{mailto:\mentoremail}{\mentoremail})}

%% De link naar je repository
\projectrepo{\repolink}

\keywords{\projectkeywords}

\begin{document}

\begin{abstract}
Door de introductie van UKG verdwijnt de rechtstreekse doorgifte van perimetertiktijden uit OnGuard, waardoor leidinggevenden het realtime inzicht in aanwezige operatoren verliezen. Dit veroorzaakt knelpunten bij ploegwissels en het tijdig opvangen van afwezigheden. Dit project levert een lichtgewicht, veilige webapplicatie in ASP.NET Core Razor Pages en C\# die actuele OnGuard-tiktijden toont en opzoekingen tot zes maanden mogelijk maakt. Toegang en autorisatie worden via Microsoft Entra ID en de Microsoft Graph API strikt afgeschermd tot de eigen directe medewerkers. De containerized applicatie draait op Kubernetes via een GitOps-flow met GitHub Actions, ArgoCD en Azure Key Vault, met gecentraliseerde monitoring via Serilog en Grafana.
\end{abstract}

%%=============================================================================
%% De inhoud van je projectvoorstel
%%=============================================================================
%%
%% Zeven rubrieken, samen hoogstens twee bladzijden. Houd het kort: per
%% rubriek volstaan één tot drie zinnen. Een voorstel is geen rapport - het
%% moet je begeleider in vijf minuten laten beslissen of dit een geschikt
%% graduaatsproject is.
%%
%% Stem je voorstel af met je mentor VOORDAT je het indient. Je mentor
%% bevestigt dat dit is wat het bedrijf bedoelt en dat de tien werkdagen
%% haalbaar zijn.

%---------- 1. Context en probleemstelling ------------------------------------

\section{Context en probleemstelling}%
\label{vst:context}

Tijdens mijn stage binnen het IT-team ter ondersteuning van HR (het People Experience team / PX) bij Volvo Car Gent werk ik aan digitale oplossingen voor personeelsbeheer en medewerkerservaring. Door de introductie van UKG verdwijnt het directe inzicht in perimetertikkingen via OnGuard aan de fabriekspoorten, omdat UKG enkel registreert aan de tikklokken bij de werkposten. Hierdoor verliezen leidinggevenden het actuele zicht op wie al op de site aanwezig is, wat leidt tot vertragingen bij shiftwissels en het tijdig opvangen van afwezige operatoren.

%---------- 2. Doelstelling ---------------------------------------------------

\section{Doelstelling}%
\label{vst:doelstelling}

Het resultaat is een beveiligde webapplicatie waarmee leidinggevenden realtime perimeteraanwezigheden per team bekijken en tiktijden tot zes maanden terug opvragen via Global ID. Het project is geslaagd zodra leidinggevenden via een responsieve interface foutloos en uitsluitend gegevens van hun eigen directe medewerkers kunnen inzien binnen een geautomatiseerde Ku\-ber\-ne\-tes-om\-ge\-ving.

%---------- 3. Technologieën --------------------------------------------------

\section{Technologieën}%
\label{vst:technologieen}

De frontend gebruikt ASP.NET Core Razor Pages (server-rendered en lichtgewicht), gekoppeld aan een C\#-backend die communiceert met de OnGuard SQL-database. Teamgegevens worden onderzocht aan de hand van Kafka Events en een REST API binnen de HR Service Layer die data uit SAP SuccessFactors ontsluit. Authenticatie verloopt via Microsoft Entra ID om Single Sign-On (SSO) te bieden, gecombineerd met Microsoft Graph API voor hiërarchische autorisatie. De deployment verloopt via GitHub Actions en ArgoCD naar Kubernetes met Azure Key Vault, aangevuld met Serilog en Grafana voor logging. Nieuw voor mij zijn Kafka, ArgoCD, Kubernetes-deployment en Microsoft Graph API. AI wordt gebruikt voor coderefactoring en pipeline-hulp. Wegens de verwerking van persoonsgegevens zijn de AVG-principes van dataminimalisatie en rolgebaseerde afscherming strikt van toepassing.

%---------- 4. Afbakening -----------------------------------------------------

\section{Afbakening}%
\label{vst:afbakening}

Mijn reguliere stagewerk bestaat uit softwareontwikkeling en onderhoud in het team. Mijn graduaatsproject is het zelfstandig ontwerpen, bouwen en deployen van de perimetertiktijden-applicatie: van de OnGuard SQL-koppeling en autorisatielogica tot de containerisatie, ArgoCD-uitrol en monitoring. Buiten scope vallen hardware-aanpassingen aan poorten of tikklokken, synchronisatie naar UKG en payroll-verwerking. Realtime updates via WebSockets gelden als uitbreiding bij tijdsoverschot.

%---------- 5. Verwachte meerwaarde -------------------------------------------

\section{Verwachte meerwaarde}%
\label{vst:meerwaarde}

De applicatie herstelt het cruciale overzicht bij ploegwissels, zodat onbemande posten direct opgevangen worden en productiestilstand voorkomen wordt. Daarnaast borgt de Microsoft Graph-autorisatie dat tiktijden AVG-conform worden afgeschermd, en levert het project een herbruikbaar GitOps-sjabloon op voor toekomstige interne microservices.

%---------- 6. Planning en werkbelasting --------------------------------------

\section{Planning en werkbelasting}%
\label{vst:planning}

Het project omvat 100 uren (tien werkdagen), gespreid in vaste projectblokken in overleg met mijn mentor.

\begin{table}[htbp]
  \centering
  \small
  \begin{tabular}{llr}
    \toprule
    \textbf{Fase} & \textbf{Wat lever je op} & \textbf{Uren} \\
    \midrule
    Analyse en vereisten      & Vereisten en OnGuard-datamapping        & 15  \\
    Technologiekeuze \& PoC   & PoC Microsoft Graph en datakoppeling    & 10  \\
    Ontwerp en architectuur   & Architectuur, queries en UI-wireframes  & 15  \\
    Realisatie                & Razor Pages UI, backend en Graph-logica & 35  \\
    DevOps \& Monitoring      & Dockerfile, CI/CD, ArgoCD en Grafana    & 10  \\
    Testen en validatie       & Rolvalidatie en tests leidinggevenden   & 5   \\
    Documentatie en demo      & Rapportage, poster en demonstratie      & 10  \\
    \midrule
    \textbf{Totaal}           &                                         & \textbf{100} \\
    \bottomrule
  \end{tabular}
  \caption[Planning en werkbelasting]{\label{tab:vst-planning}%
    Planning van het graduaatsproject (100 uren).}
\end{table}

%---------- 7. Risico's -------------------------------------------------------

\section{Risico's}%
\label{vst:risicos}

\begin{itemize}
    \item \textbf{Beveiligde toegang tot OnGuard en SAP SuccessFactors:} Rechtenaanvragen via het \emph{least-privilege}-principe kunnen doorlooptijd vragen. \emph{Terugvalplan:} Ontwikkelen met geanonimiseerde mock-data en views.
    \item \textbf{Feedback en businessafstemming:} Bruikbaarheid vereist directe input van de werkvloer. \emph{Terugvalplan:} Vroegtijdige wireframes en korte reviewmomenten met mentor en sleutelgebruikers.
    \item \textbf{Beperkingen op tenant-rechten in Microsoft Graph:} Indien rechten op \emph{direct reports} te restrictief blijken in productie. \emph{Terugvalplan:} Uitwijken naar Entra ID-beveiligingsgroepen of de HR Service Layer.
\end{itemize}

\printbibliography[heading=bibintoc]

\end{document}
